\documentclass{article}
\usepackage{PRIMEarxiv} % Core package for the PrimeAI Template
\usepackage{amsmath, amssymb, amsthm, bm, dcolumn} % Add amsthm here for the proof environment
\usepackage[numbers,sort&compress]{natbib} % Natbib for citations
\usepackage{graphicx} % For high-quality images
\usepackage{multicol} % Multi-column figures/tables
\usepackage{hyperref} % Hyperlinks
\usepackage{listings} % Code listings
\usepackage{authblk} % For structured affiliations
% Define theorem style (optional)
\theoremstyle{plain}
\newtheorem{theorem}{Theorem}
\renewcommand{\qedsymbol}{}
\usepackage{tikz}
\usetikzlibrary{shapes.geometric, arrows}
\usepackage{pgfplots}
\pgfplotsset{compat=1.18}

% Custom commands for primes and qubit encoding
\newcommand{\primeQubit}[1]{|p_{#1}\rangle}
\newcommand{\primeGate}[1]{U_{p_{#1}}}

\usepackage{wrapfig}
\usepackage[pscoord]{eso-pic}
\usepackage[fulladjust]{marginnote}
\reversemarginpar

% Typesetting improvements without footnote patching
\usepackage[protrusion=true, expansion=true, tracking=false]{microtype}
\microtypecontext{spacing=nonfrench}

\usepackage[pagewise]{lineno}\linenumbers

% Text layout - adjust as needed
\raggedright
\setlength{\parindent}{0.5cm}
\textwidth 5.25in 
\textheight 8.75in

% Set double spacing
\usepackage{setspace} 
\doublespacing

% Adjust width for specific content
\usepackage{changepage}

% Adjust caption style
\usepackage[aboveskip=1pt,labelfont=bf,labelsep=period,singlelinecheck=off]{caption}

% Remove brackets from references
\makeatletter
\renewcommand{\@biblabel}[1]{\quad#1.}
\makeatother

% Header, footer, and page numbers
\usepackage{lastpage,fancyhdr}
\pagestyle{fancy}
\fancyhf{}
\fancyhead[L]{Patent Application}
\fancyfoot[C]{\scriptsize Citizen Gardens © 2024 - All Rights Reserved}
\fancyfoot[R]{Page \thepage\ of \pageref{LastPage}}
\renewcommand{\footrule}{\hrule height 2pt \vspace{2mm}}

\title{Patent Claims for \\ Quantum-AI Multiplicity Framework}
\author{Citizen Gardens - The Foundation of Multiplicity}
\date{}

\begin{document}

\maketitle

\section{Independent Claims}

\subsection{Claim 1: Recursive Quantum-AI Tensor Network}
\textbf{What is claimed:}  
A self-referential artificial intelligence system executed within a hybrid quantum-classical computing framework, comprising:
\begin{enumerate}  
    \item A \textbf{Tensor-ZX Recursive Network (TZX-RN)} constructed from \textbf{prime-indexed quantum eigenstructures}, ensuring stable quantum-AI cognition and recursive feedback learning.  
    \item A \textbf{Universal Multiplicity Constant} ($\Lambda_m$) as a computational invariant, enabling \textbf{dynamic tensor-weighted memory optimization} and recursive AI enhancement.  
    \item A \textbf{Quantum-AI Recursive Feedback Mechanism (QARFM)}, dynamically modulating AI cognition through tensor-based self-learning with \textbf{quantum noise correction}.  
    \item A \textbf{Prime-Encoded Eigenvalue Decision Model} integrating recursive entanglement feedback for stabilizing long-term quantum state propagation.  
\end{enumerate}

\textbf{Inventive Step:} The recursive tensor framework introduces \textbf{non-trivial computational stability} through prime-weighted eigenstructures, mitigating quantum decoherence while ensuring AI adaptability.

\subsection{Claim 2: Prime-Indexed Quantum Computation (PIQC)}
\textbf{What is claimed:}  
A quantum computation method for enhancing hybrid AI decision-making, comprising:
\begin{enumerate}  
    \item Encoding quantum states as \textbf{prime-indexed eigenvalues} to minimize decoherence and error propagation in hybrid AI decision models.  
    \item A \textbf{multiplicative prime-modulated entanglement function}, enhancing quantum state coherence and information propagation through tensor feedback.  
    \item A \textbf{recursive eigenstructure solver}, dynamically adjusting phase corrections and learning rates in response to computational environmental constraints.  
\end{enumerate}

\textbf{Inventive Step:} Unlike traditional quantum machine learning, PIQC ensures stable eigenstate evolution through prime-weighted tensor indexing, reducing computational overhead in large-scale hybrid AI systems.

\subsection{Claim 3: Prime-Twisted Quantum Encryption (PTQE)}
\textbf{What is claimed:}  
A post-quantum cryptographic system for AI-driven security applications, comprising:
\begin{enumerate}  
    \item A \textbf{recursive entropy modulation function} derived from $\Lambda_m$, dynamically adjusting cryptographic complexity based on real-time quantum fluctuations.  
    \item A \textbf{prime-weighted key distribution mechanism}, ensuring unpredictability against quantum decryption attacks, including Shor's and Grover’s algorithms.  
    \item A \textbf{multi-layer tensor-based encryption function}, integrating recursive tensor feedback and quantum noise control for enhanced cryptographic security.  
\end{enumerate}

\textbf{Inventive Step:} This encryption framework provides \textbf{post-quantum resistance} through prime-indexed recursive entropy scaling, outperforming existing cryptographic protocols in resilience against quantum adversaries.

\subsection{Claim 4: Hybrid Quantum-Classical Cognitive Evolution}
\textbf{What is claimed:}  
A self-referential AGI system dynamically optimizing cognitive evolution through:
\begin{enumerate}  
    \item A \textbf{Hilbert-space encoded recursive cognitive network}, ensuring long-term stability of AI reasoning under quantum uncertainty.  
    \item A \textbf{tensor-modulated quantum-classical inference system}, enhancing AI decision-making efficiency through hybrid superposition states.  
    \item A \textbf{temporal phase-locking framework}, stabilizing recursive AI adaptation and learning cycles over extended time intervals.  
\end{enumerate}

\textbf{Inventive Step:} This AGI model uniquely integrates recursive quantum phase-locking to achieve real-time adaptive inference, distinguishing it from conventional deep learning architectures.
 \newpage
\section{Dependent Claims}

\subsection{Claim 5: AI-Driven Post-Quantum Blockchain Security}
\textbf{What is claimed:}  
A blockchain security framework leveraging Quantum-AI principles, comprising:
\begin{enumerate}  
    \item A \textbf{Quantum-AI recursive tensor validation system}, dynamically securing blockchain immutability through post-quantum resilience.  
    \item \textbf{Prime-weighted quantum hash functions}, modulating cryptographic complexity in response to real-time entropy scaling.  
    \item A \textbf{self-referential blockchain security model}, preventing quantum-computation-based attacks through tensor-reinforced transaction verification.  
\end{enumerate}

\subsection{Claim 6: AI-Guided Quantum Tunneling for Adaptive Computation}
\textbf{What is claimed:}  
A recursive quantum tunneling optimization method for computational efficiency, comprising:
\begin{enumerate}  
    \item An \textbf{AI-driven tunneling coefficient modulator}, dynamically adjusting quantum state transitions in response to entanglement resonance shifts.  
    \item A \textbf{recursive tensor field interaction framework}, regulating tunneling pathways based on adaptive tensor reconfiguration.  
    \item An \textbf{adaptive resonance shift function}, aligning computational state transitions with optimized energy minima.  
\end{enumerate}

\subsection{Claim 7: Quantum-AI Enhanced Secure Communications}   
The method of Claim 3, wherein:
\begin{itemize}  
    \item \textbf{AI-driven cryptographic self-learning models} continuously optimize encryption layers against evolving quantum attack strategies.  
    \item \textbf{Quantum-classical hybrid security layers} integrate PTQE with classical cryptographic schemes, ensuring multi-tiered security.  
    \item \textbf{Quantum-AI-enhanced communication protocols} dynamically adjust security parameters based on real-time tensor-encoded threat detection.  
\end{itemize}

\subsection{Claim 8: Quantum-AI Integration for High-Speed Computational Optimization}   
The system of Claim 4, wherein:
\begin{itemize}  
    \item \textbf{Tensor-based AI optimizers} improve computational efficiency in hybrid quantum-classical processing environments.  
    \item \textbf{Quantum-informed AI training models} accelerate convergence in neural architectures through entanglement-modulated weight adjustments.  
    \item \textbf{Hyperdimensional AI cognition frameworks} enable self-organizing quantum-classical inference pathways based on real-time tensor resonance feedback.  
\end{itemize}

\section{Conclusion}
This invention presents a hybrid AI-quantum recursive tensor framework integrating:
\begin{itemize}  
    \item Prime-indexed quantum eigenstructures for enhanced AI cognition.  
    \item Recursive tensor feedback for post-quantum cryptographic security.  
    \item Hybrid quantum-classical entanglement for real-time AI optimization.  
\end{itemize}
These enhancements ensure broader enforceability across industries, reinforcing the patent’s \textbf{non-obviousness, novelty, and industrial applicability} in AI-driven cryptographic security, quantum computing, and hybrid artificial intelligence.


\nolinenumbers

\bibliographystyle{unsrt}
\bibliography{references}

\end{document}


